% This configuration is adapted from: https://github.com/sleepymalc/LaTeX-Template/blob/main/Note/header.tex

%----------------------------------------------------------------------------------------
%	FORMATTING
%----------------------------------------------------------------------------------------

\setlength{\parindent}{2em}
\setlength{\parskip}{0.5\baselineskip}
\setlength{\voffset}{-15pt}

%----------------------------------------------------------------------------------------
%	PACKAGES AND OTHER DOCUMENT CONFIGURATIONS
%----------------------------------------------------------------------------------------

\usepackage[a4paper, top=2cm, bottom=2cm, left=2.5cm, right=2.5cm]{geometry} % Sets margin to 2.5cm for A4 Paper
\usepackage[onehalfspacing]{setspace} % Sets Spacing to 1.5

\usepackage[T1]{fontenc} % Use European encoding
\usepackage[utf8]{inputenc} % Use UTF-8 encoding
%\usepackage{CJKutf8} % Chinese
\usepackage{charter} % Use the Charter font cuz why not
%\usepackage{newpxtext} % Use the NewPX font cuz why not
%\usepackage{newpxmath} % Use the NewPX font cuz why not
\usepackage{microtype} % Slightly tweak font spacing for aesthetics

\usepackage[english]{babel} % Language hyphenation and typographical rules

\usepackage{amsthm, amsmath, amssymb} % Mathematical typesetting
\numberwithin{equation}{section}
\newcommand\numberthis{\addtocounter{equation}{1}\tag{\theequation}} % number this equation

\usepackage{marvosym, wasysym} % More symbols
\usepackage{centernot} % centered not
\usepackage{float} % Improved interface for floating objects
\usepackage[final, colorlinks = true, 
            linkcolor = cyan, 
            citecolor = violet,
            urlcolor = magenta]{hyperref} % For hyperlinks in the PDF
\usepackage{graphicx, multicol} % Enhanced support for graphics
\usepackage[usenames,dvipsnames]{xcolor} % Driver-independent color extensions
\usepackage{rotating} % Rotation tools
\usepackage{listings, style/lstlisting} % Environment for non-formatted code, !uses style file!
\usepackage{pseudocode} % Environment for specifying algorithms in a natural way
\usepackage{style/avm} % Environment for f-structures, !uses style file!
\usepackage{booktabs} % Enhances quality of tables

\usepackage{tikz-qtree} % Easy tree drawing tool
\usetikzlibrary{arrows.meta,decorations.pathmorphing,backgrounds,positioning,fit,petri,calc}
\tikzset{every tree node/.style={align=center,anchor=north},
         level distance=2cm} % Configuration for q-trees
\usepackage{style/btree} % Configuration for b-trees and b+-trees, !uses style file!

\usepackage{titlesec} % Allows customization of titles
\renewcommand\thesection{\arabic{section}} % Arabic numerals for the sections
\titleformat{\section}{\Large\bfseries}{\thesection}{1em}{}
\renewcommand\thesubsection{\alph{subsection})} % Alphabetic numerals for subsections
\titleformat{\subsection}{\large\bfseries}{\thesubsection}{1em}{}
\renewcommand\thesubsubsection{\roman{subsubsection}.} % Roman numbering for subsubsections
\titleformat{\subsubsection}{\large}{\thesubsubsection}{1em}{}

\usepackage[all]{nowidow} % Removes widows

\usepackage{framed} % framed

%\usepackage[backend=biber,style=alphabetic,
%            sorting=nyt, natbib=true]{biblatex} % Complete reimplementation of bibliographic facilities


\usepackage{csquotes} % Context sensitive quotation facilities

\usepackage[yyyymmdd]{datetime} % Uses YEAR-MONTH-DAY format for dates
\renewcommand{\dateseparator}{-} % Sets dateseparator to '-'

\usepackage{fancyhdr} % Headers and footers
\pagestyle{fancy} % All pages have headers and footers
\fancyhead{}\renewcommand{\headrulewidth}{0pt} % Blank out the default header
\fancyfoot[L]{} % Custom footer text
\fancyfoot[C]{} % Custom footer text
\fancyfoot[R]{} % Custom footer text

\usepackage{lipsum}
\usepackage{enumitem}
\usepackage{caption}
\captionsetup{belowskip=0pt}

\usepackage{bigdelim} % for the big braces

\usepackage{makecell} % for cool table cells \makecell




% correct
\definecolor{correct}{HTML}{009900}
\newcommand\correct[2]{\ensuremath{\:}{\color{red}{#1}}\ensuremath{\to }{\color{correct}{#2}}\ensuremath{\:}}
\newcommand\green[1]{{\color{correct}{#1}}}

% hide parts
\newcommand\hide[1]{}

% footnote
\newcommand\blfootnote[1]{%
  \begingroup
  \renewcommand\thefootnote{}\footnote{#1}%
  \addtocounter{footnote}{-1}%
  \endgroup
}

% Enables comments in red on margin
\newcommand{\comment}[1]{\marginpar{\scriptsize \textcolor{red}{#1}}} 

% Clear out the item autoref name
\makeatletter
\def\itemautorefname{\@gobble}
\makeatother

% set "enumerate" env spacing
\setlist[enumerate]{topsep=0pt, partopsep=0pt}

% invisible sections (copied from https://tex.stackexchange.com/a/68296)
\newcommand\invisiblesection[1]{%
  \refstepcounter{section}%
  \addcontentsline{toc}{section}{\protect\numberline{\thesection}#1}%
  \sectionmark{#1}}

%----------------------------------------------------------------------------------------
%	MATH CONFIGURATIONS
%----------------------------------------------------------------------------------------
\usepackage{mathtools}
\usepackage{fbox}
%\usepackage{etoolbox}
\usepackage{twemojis} % for emojis
\usepackage{esvect} % cooler vectors \vv

% =========================
% elementry operations
\newcommand{\gt}{\mathrel{>}}
\newcommand{\lt}{\mathrel{<}}
\newcommand{\equals}{\mathrel{=}}
\DeclareMathOperator{\sgn}{sgn}
\let\oldemptyset\emptyset
\let\emptyset\varnothing
\apptocmd{\lim}{\limits}{}{}
%\let\implies\Rightarrow
%\let\impliedby\Leftarrow
%\let\iff\Leftrightarrow
\let\epsilon\varepsilon
\newcommand{\minus}{\scalebox{0.75}[1.0]{$-$}}
\renewcommand{\bar}[1]{\overline{#1}}
\renewcommand{\tilde}[1]{\widetilde{#1}}
\DeclarePairedDelimiter\abs{\lvert}{\rvert}
\DeclarePairedDelimiter\norm{\lVert}{\rVert}
\newcommand{\concat}{\mathbin{\Vert}}
\newcommand{\bit}{\{0,1\}}
\newcommand{\powerset}[1]{\mathcal{P}(#1)}
\newcommand{\symdiff}{\mathop{\triangle}}
\newcommand{\func}{\mathsf{Func}}
\newcommand{\fst}{\mathop{\textrm{fst}}}
\newcommand{\snd}{\mathop{\textrm{snd}}}

% =========================
% logical operations
\newcommand{\entails}[1]{\vdash_{\mathsf{#1}}}

% =========================
% probabilistic operations
\DeclareMathOperator*{\Var}{Var}
\DeclareMathOperator*{\Exp}{\mathbb{E}}
\newcommand{\ensemble}[2]{\{#1\}_{#2\in\Nbb}}
\newcommand{\sample}{\mathrel{\leftarrow}}
\newcommand{\Usample}{\mathrel{\xleftarrow{\$}}}
\newcommand{\prdist}[2]{\pr\limits_{#1}\left[\ #2\ \right]}
\newcommand{\iid}{\overset{\mathrm{iid}}{\sim}}
\newcommand{\eqdist}{\overset{\mathrm{d}}{=}}
\newcommand{\indep}{\mathrel{\perp\!\!\!\perp}}

% =========================
% cryotographic notations
% 0. basic defs
\newcommand{\ind}{\approx} % indistinguishable
\newcommand{\aympind}{\approx_c} % asymptotically indistinguishable
\newcommand{\adv}{\mathop{\mathrm{Adv}}}
% 1. Primitive / scheme names
\newcommand{\PKE}{\mathsf{PKE}}
\newcommand{\SKE}{\mathsf{SKE}}
\newcommand{\KEM}{\mathsf{KEM}}
\newcommand{\OWPKE}{\mathsf{PKE}_{\mathsf{OW}}}
\newcommand{\CCAPKE}{\mathsf{PKE}_{\mathsf{CCA}}}
\newcommand{\LWE}{\mathsf{LWE}}
\newcommand{\DLWE}{\mathsf{DLWE}}
\newcommand{\SLWE}{\mathsf{SLWE}}
% 2. Schemes
\newcommand{\Gen}{\mathsf{KeyGen}}
\newcommand{\Setup}{\mathsf{Setup}}
\newcommand{\Enc}{\mathsf{Enc}}
\newcommand{\Dec}{\mathsf{Dec}}
\newcommand{\Encaps}{\mathsf{Encaps}}
\newcommand{\Decaps}{\mathsf{Decaps}}
\newcommand{\Sign}{\mathsf{Sign}}
\newcommand{\Vrfy}{\mathsf{Vrfy}}
\newcommand{\Com}{\mathsf{Com}}
\newcommand{\Open}{\mathsf{Open}}
\newcommand{\view}{\mathsf{view}}
\newcommand{\out}{\mathsf{out}}
\newcommand{\Ext}{\mathsf{Ext}}
\newcommand{\Sim}{\mathsf{Sim}}
\newcommand{\Trans}{\mathsf{Trans}}
\newcommand{\RO}{\mathsf{RO}}
% 3. Keys / ciphertexts / common objects
\newcommand{\pk}{\mathsf{pk}}
\newcommand{\sk}{\mathsf{sk}}
\newcommand{\ct}{\mathsf{ct}}
\newcommand{\crs}{\mathsf{crs}}
\newcommand{\st}{\mathsf{st}}
\newcommand{\accept}{\mathsf{accept}}
\newcommand{\reject}{\mathsf{reject}}
% 4. Functions
\newcommand{\negl}{\mathsf{negl}}
\newcommand{\nonnegl}{\mathsf{nonnegl}}
\newcommand{\poly}{\mathsf{poly}}
% 5. Entities in games / proofs
\newcommand{\Adv}{\mathbf{Adv}}
\newcommand{\Ch}{\mathsf{Ch}}
\newcommand{\Hyb}{\mathsf{Hyb}}
\renewcommand{\Game}{\mathbf{\mathsf{G}}}
% 6. Protocol-specific notation
\newcommand{\com}{\mathsf{com}}
\newcommand{\ch}{\mathsf{ch}}
\newcommand{\resp}{\mathsf{resp}}
\newcommand{\commit}{\mathsf{commit}}
\newcommand{\interact}[1]{\left<#1\right>}

% =========================
% shorthands for convenience
\newcommand{\Ccal}{\mathcal{C}}
\newcommand{\Acal}{\mathcal{A}}
\newcommand{\Rcal}{\mathcal{R}}
\newcommand{\Mcal}{\mathcal{M}}
\newcommand{\Kcal}{\mathcal{K}}
\newcommand{\Ical}{\mathcal{I}}
\newcommand{\Tcal}{\mathcal{T}}
\newcommand{\Dcal}{\mathcal{D}}
\newcommand{\Gcal}{\mathcal{G}}
\newcommand{\Ocal}{\mathcal{O}}
\newcommand{\Hcal}{\mathcal{H}}
\newcommand{\Xcal}{\mathcal{X}}
\newcommand{\Ycal}{\mathcal{Y}}
\newcommand{\Zcal}{\mathcal{Z}}
\newcommand{\Nbb}{\mathbb{N}}
\newcommand{\Zbb}{\mathbb{Z}}
\newcommand{\Rbb}{\mathbb{R}}
\newcommand{\Gbb}{\mathbb{G}}
\newcommand{\Tsf}{\mathsf{T}}
\newcommand{\Usf}{\mathsf{U}}
\newcommand{\Lfrak}{\mathfrak{L}}

% =========================
% arrows for sequence diagram
\def\XLR#1{\xleftrightarrow[\rule{2cm}{0pt}]{#1}}% \rule defines the width
\def\XR#1{\xrightarrow[\rule{2cm}{0pt}]{#1}}
\def\XL#1{\xleftarrow[\rule{2cm}{0pt}]{#1}}

\newcommand{\xdownarrow}[2][0.5cm]{%
  \mathrel{\underset{#2}{\left\downarrow\vbox to #1{}\right.\kern-\nulldelimiterspace}}%
}
\newcommand{\XD}[1]{\xdownarrow[0.5cm]{#1}} % \rule defines the height

% =========================
% misc
\newcommand{\highlight}[2]{\colorbox{#1}{$\displaystyle #2$}}
\newcommand{\proofcase}[1]{%
  \par\medskip
  \noindent\textbf{\large #1.}\enspace
  \par\smallskip
}
%\newcommand{\too}{\mathbin{\Rightarrow}}
%\newcommand\tildee{\raise.17ex\hbox{$\scriptstyle\sim$}}
% \newcommand{\sus}{
%     \textrm{\fontspec{Noto Sans Sinhala}ඞ}
% }


%----------------------------------------------------------------------------------------
%	THEOREM CONFIGURATIONS
%----------------------------------------------------------------------------------------

% \newtheorem{theorem}{Theorem}[section]
% \newtheorem{assumption}{Assumption}
% \newtheorem{proposition}[theorem]{Proposition}
% \newtheorem{conjecture}[theorem]{Conjecture}
% \newtheorem{lemma}[theorem]{Lemma}
% \newtheorem{claim}[theorem]{Claim}
% \newtheorem{fact}[theorem]{Fact}
% \newtheorem{corollary}[theorem]{Corollary}
% \newtheorem*{theorem*}{Theorem}
% \newtheorem*{assumption*}{Assumption}

% \theoremstyle{remark}
% \newtheorem{remark}[theorem]{Remark}

% \theoremstyle{definition}
% \newtheorem{definition}[theorem]{Definition}
% \newtheorem*{definition*}{Definition}
% \newtheorem{example}[theorem]{Example}
% \newtheorem{oq}[theorem]{Open Question}


\makeatother
\usepackage{thmtools}
\usepackage[framemethod=TikZ]{mdframed}

\mdfsetup{skipabove=1em,skipbelow=0em}

\theoremstyle{definition}

% simple textbox
\mdfdefinestyle{simplebox}{
    linewidth=1pt,
    linecolor=black,
    backgroundcolor=white,
    leftline=true,
    rightline=true,
    topline=true,
    bottomline=true,
    roundcorner=0pt,
    skipabove=0.3em,
    skipbelow=0em
}



\declaretheoremstyle[
	headfont=\bfseries\sffamily\color{ForestGreen!70!black}, bodyfont=\normalfont,
	mdframed={
			linewidth=2pt,
			rightline=false, topline=false, bottomline=false,
			linecolor=ForestGreen, backgroundcolor=ForestGreen!5,
			nobreak=false
		}
]{thmgreenbox}

\declaretheoremstyle[
	headfont=\bfseries\sffamily\color{ForestGreen!70!black}, bodyfont=\normalfont,
	mdframed={
			linewidth=2pt,
			rightline=false, topline=false, bottomline=false,
			linecolor=ForestGreen, backgroundcolor=ForestGreen!8,
			nobreak=false
		}
]{thmgreen2box}

\declaretheoremstyle[
	headfont=\bfseries\sffamily\color{ForestGreen!70!black}, bodyfont=\normalfont,
	mdframed={
			linewidth=2pt,
			rightline=false, topline=false, bottomline=false,
			linecolor=ForestGreen,
			nobreak=false
		}
]{thmgreenline}

\declaretheoremstyle[
	headfont=\bfseries\sffamily\color{NavyBlue!70!black}, bodyfont=\normalfont,
	mdframed={
			linewidth=2pt,
			rightline=false, topline=false, bottomline=false,
			linecolor=NavyBlue, backgroundcolor=NavyBlue!5,
			nobreak=false
		}
]{thmbluebox}

\declaretheoremstyle[
	headfont=\bfseries\sffamily\color{TealBlue!70!black}, bodyfont=\normalfont,
	mdframed={
			linewidth=2pt,
			rightline=false, topline=false, bottomline=false,
			linecolor=TealBlue,
			nobreak=false
		}
]{thmblueline}

\declaretheoremstyle[
	headfont=\bfseries\sffamily\color{RawSienna!70!black}, bodyfont=\normalfont,
	mdframed={
			linewidth=2pt,
			rightline=false, topline=false, bottomline=false,
			linecolor=RawSienna, backgroundcolor=RawSienna!5,
			nobreak=false
		}
]{thmredbox}

\declaretheoremstyle[
	headfont=\bfseries\sffamily\color{RawSienna!70!black}, bodyfont=\normalfont,
	mdframed={
			linewidth=2pt,
			rightline=false, topline=false, bottomline=false,
			linecolor=RawSienna, backgroundcolor=RawSienna!8,
			nobreak=false
		}
]{thmred2box}

\declaretheoremstyle[
	headfont=\bfseries\sffamily\color{SeaGreen!70!black}, bodyfont=\normalfont,
	mdframed={
			linewidth=2pt,
			rightline=false, topline=false, bottomline=false,
			linecolor=SeaGreen, backgroundcolor=SeaGreen!2,
			nobreak=false
		}
]{thmgreen3box}

\declaretheoremstyle[
	headfont=\bfseries\sffamily\color{WildStrawberry!70!black}, bodyfont=\normalfont,
	mdframed={
			linewidth=2pt,
			rightline=false, topline=false, bottomline=false,
			linecolor=WildStrawberry, backgroundcolor=WildStrawberry!5,
			nobreak=false
		}
]{thmpinkbox}

\declaretheoremstyle[
	headfont=\bfseries\sffamily\color{WildStrawberry!70!black}, bodyfont=\normalfont,
	mdframed={
			linewidth=2pt,
			rightline=false, topline=false, bottomline=false,
			linecolor=WildStrawberry,
			nobreak=false
		}
]{thmpinkline}

\declaretheoremstyle[
	headfont=\bfseries\sffamily\color{MidnightBlue!70!black}, bodyfont=\normalfont,
	mdframed={
			linewidth=2pt,
			rightline=false, topline=false, bottomline=false,
			linecolor=MidnightBlue, backgroundcolor=MidnightBlue!5,
			nobreak=false
		}
]{thmblue2box}

\declaretheoremstyle[
	headfont=\bfseries\sffamily\color{Gray!70!black}, bodyfont=\normalfont,
	mdframed={
			linewidth=2pt,
			rightline=false, topline=false, bottomline=false,
			linecolor=Gray, backgroundcolor=Gray!5,
			nobreak=false
		}
]{notgraybox}

\declaretheoremstyle[
	headfont=\bfseries\sffamily\color{Gray!70!black}, bodyfont=\normalfont,
	mdframed={
			linewidth=2pt,
			rightline=false, topline=false, bottomline=false,
			linecolor=Gray,
			nobreak=false
		}
]{notgrayline}

% \declaretheoremstyle[
% 	headfont=\bfseries\sffamily\color{RawSienna!70!black}, bodyfont=\normalfont,
% 	numbered=no,
% 	mdframed={
% 			linewidth=2pt,
% 			rightline=false, topline=false, bottomline=false,
% 			linecolor=RawSienna, backgroundcolor=RawSienna!1,
% 		},
% 	qed=\qedsymbol
% ]{thmproofbox}

\declaretheoremstyle[
	headfont=\bfseries\sffamily\color{NavyBlue!70!black}, bodyfont=\normalfont,
	numbered=no,
	mdframed={
			linewidth=2pt,
			rightline=false, topline=false, bottomline=false,
			linecolor=NavyBlue, backgroundcolor=NavyBlue!1,
			nobreak=false
		}
]{thmexplanationbox}

\declaretheoremstyle[
	headfont=\bfseries\sffamily\color{WildStrawberry!70!black}, bodyfont=\normalfont,
	numbered=no,
	mdframed={
			linewidth=2pt,
			rightline=false, topline=false, bottomline=false,
			linecolor=WildStrawberry, backgroundcolor=WildStrawberry!1,
			nobreak=false
		}
]{thmanswerbox}

\declaretheoremstyle[
	headfont=\bfseries\sffamily\color{Violet!70!black}, bodyfont=\normalfont,
	mdframed={
			linewidth=2pt,
			rightline=false, topline=false, bottomline=false,
			linecolor=Violet, backgroundcolor=Violet!1,
			nobreak=false
		}
]{conjpurplebox}

% shared counter "stmt", resets within section
\declaretheorem[style=thmblue2box, name=, numberwithin=section]{stmt}

\declaretheorem[style=thmgreenbox, name=Definition, numberwithin=section]{definition}
\declaretheorem[style=thmgreen2box, name=Definition, numbered=no]{definition*}
\declaretheorem[style=thmredbox, name=Theorem, numberwithin=section]{theorem}
\declaretheorem[style=thmred2box, name=Theorem, numbered=no]{theorem*}
\declaretheorem[style=thmredbox, name=Lemma, numberwithin=section]{lemma}
\declaretheorem[style=thmred2box, name=Lemma, numbered=no]{lemma*}
\declaretheorem[style=thmredbox, name=Proposition, numberwithin=section]{proposition}
\declaretheorem[style=thmredbox, name=Corollary, numberwithin=section]{corollary}
\declaretheorem[style=thmpinkbox, name=Problem, numberwithin=section]{problem}
\declaretheorem[style=thmpinkbox, name=Problem, numbered=no]{problem*}
\declaretheorem[style=thmblue2box, name=Claim, sibling=stmt]{claim}
\declaretheorem[style=thmblue2box, name=Claim, numbered=no]{claim*}
\declaretheorem[style=conjpurplebox, name=Conjecture, numberwithin=section]{conjecture}
\declaretheorem[style=thmblueline, name=Observation, sibling=stmt]{observation}
\declaretheorem[style=thmblueline, name=Observation, numbered=no]{observation*}

\renewcommand\theproblem{\thesection} %\renewcommand\theproblem{\thesection.\arabic {problem}}
\renewcommand\thelemma{\thesection.\arabic {lemma}}
\renewcommand\theclaim{\thesection.\arabic {claim}}
\renewcommand\theobservation{\thesection.\arabic {observation}}

\renewcommand\theHdefinition{\thesection.\arabic {definition}}
\renewcommand\theHtheorem{\thesection.\arabic {theorem}}
\renewcommand\theHlemma{\thesection.\arabic {lemma}}
\renewcommand\theHclaim{\thesection.\arabic {claim}}
\renewcommand\theHobservation{\thesection.\arabic {observation}}
\renewcommand\theHproposition{\thesection.\arabic {proposition}}
\renewcommand\theHcorollary{\thesection.\arabic {corollary}}
\renewcommand\theHproblem{\thesection {problem}} %\renewcommand\theHproblem{\thesection.\arabic {problem}}
\renewcommand\theHconjecture{\thesection.\arabic {conjecture}}

% Redefine proof environment to get a full control.
\makeatletter
\renewenvironment{proof}[1][\proofname]{\par
	\pushQED{\qed}%
	\normalfont \topsep-2\p@\@plus6\p@\relax
	\trivlist
	\item[\hskip\labelsep
	            \color{RawSienna!70!black}\sffamily\bfseries
	            #1\@addpunct{.}]\ignorespaces
	\begin{mdframed}[linewidth=2pt,rightline=false, topline=false, bottomline=false,linecolor=RawSienna, backgroundcolor=RawSienna!1]
        \setlength{\parindent}{1.5em}  % set indent here
		}{%
		\popQED\endtrivlist\@endpefalse
	\end{mdframed}
}
\makeatother

\declaretheorem[style=thmbluebox, numbered=no, name=Example]{eg}
\declaretheorem[style=thmexplanationbox, numbered=no, name=Proof]{tmpexplanation}
\newenvironment{explanation}[1][]{\vspace{-10pt}\pushQED{\(\circledast\)}\begin{tmpexplanation}}{\null\hfill\popQED\end{tmpexplanation}}

\declaretheorem[style=thmblueline, numbered=no, name=Remark]{remark}
\declaretheorem[style=thmblueline, numbered=no, name=Note]{note}
\declaretheorem[style=thmpinkbox, numbered=no, name=Exercise]{exercise}
\declaretheorem[style=notgrayline, numbered=no, name=As previously seen]{prev}
\declaretheorem[style=thmgreen3box, numbered=no, name=Intuition]{intuition}
\declaretheorem[style=thmpinkline, numbered=no, name=Idea]{idea}
\declaretheorem[style=thmgreenline, numbered=no, name=Goal]{goal}
\declaretheorem[style=notgraybox, numbered=no, name=Notation]{notation}
\declaretheorem[style=thmanswerbox, numbered=no, name=Answer]{tmpanswer}
\newenvironment{answer}[1][]{%
  \vspace{-10pt}%
  \pushQED{\(\circledast\)}%
  \ifstrempty{#1}
    {\begin{tmpanswer}}
    {\begin{tmpanswer}[#1]}%
}{%
  \null\hfill\popQED\end{tmpanswer}%
}

% simple textbox
\newenvironment{simplebox}
{\begin{mdframed}[style=simplebox]}
{\end{mdframed}}


\usepackage{etoolbox}
\renewcommand{\qed}{\null\hfill\(\blacksquare\)}

\makeatletter


\usepackage{comment}
\excludecomment{wrap}

%%% https://tex.stackexchange.com/a/439785
\newcounter{relctr} %% <- counter for relations
\everydisplay\expandafter{\the\everydisplay\setcounter{relctr}{0}} %% <- reset every eq
\renewcommand*\therelctr{\alph{relctr}} %% <- label format

\newcommand\labelrel[2]{%
  \begingroup
    \refstepcounter{relctr}%
    \stackrel{\textnormal{(\alph{relctr})}}{\mathstrut{#1}}%
    \originallabel{#2}%
  \endgroup
}
\AtBeginDocument{\let\originallabel\label} %% <- store original definition


%----------------------------------------------------------------------------------------
%	MISC
%----------------------------------------------------------------------------------------
% Fix some stuff
% %http://tex.stackexchange.com/questions/76273/multiple-pdfs-with-page-group-included-in-a-single-page-warning
\pdfsuppresswarningpagegroup=1